\documentclass[11pt]{article}
\usepackage{amsmath,amssymb,amsthm,mathtools}
\usepackage[margin=1in]{geometry}
\usepackage{hyperref}

\newtheorem{theorem}{Theorem}[section]
\newtheorem{lemma}[theorem]{Lemma}
\newtheorem{proposition}[theorem]{Proposition}
\newtheorem{corollary}[theorem]{Corollary}
\theoremstyle{definition}
\newtheorem{remark}[theorem]{Remark}
\newtheorem{definition}[theorem]{Definition}
\newtheorem{notation}[theorem]{Notation}

\DeclareMathOperator{\Lam}{\Lambda}
\newcommand{\Z}{\mathbb{Z}}
\newcommand{\Hop}{H}
\newcommand{\iotain}{\iota_{12}}
\newcommand{\iotaout}{\iota_{21}}

\title{The two-parameter Hall--Littlewood exchange relation,\\
and what actually lives on the hyperbola $st=1$}
\author{Clio}
\date{8 September 2026}

\begin{document}
\maketitle

\begin{center}\small
\fbox{\parbox{0.92\textwidth}{\centering
\textbf{Author:} Clio \quad\textbullet\quad \textbf{Date:} 8 September 2026 \quad\textbullet\quad
\textbf{Recipient:} Rick (cc Robin)\\
\textbf{Title:} The two-parameter Hall--Littlewood exchange relation, and what actually lives
on the hyperbola $st=1$\\
\textbf{Corresponds to:} \texttt{clio-vega/proofs@0064531} (paper and registry nodes)\\
\textbf{Reviewed:} self-review of 8 September 2026,
\texttt{clio-vega/rick-review@e2f58ff} --- Theorems~\ref{thm:A} and \ref{thm:B} independently
re-derived ($252/252$ and $288/288$) on an h-basis engine sharing no code path with \S4's;
\S5 and the citation sentence of \S2 corrected as a result (findings F1 and F3).
}}
\end{center}
\bigskip

\begin{abstract}
Let $\Hop_t(z)$ be the Hall--Littlewood vertex operator, $\Hop_t(z)f[X]=\sigma[Xz]\,f[X-(1-t)/z]$.
We prove that for two \emph{independent} parameters $s,t$,
\[
  (z_1-t z_2)\,\Hop_t(z_1)\Hop_s(z_2)\;=\;(s z_1-z_2)\,\Hop_s(z_2)\Hop_t(z_1),
\]
by showing that both sides equal $(z_1-z_2)\Phi$ for one and the same normal-ordered
series $\Phi$. Specialising $s=t$ recovers Zabrocki's relation (2.14). The exchange factor
$(sz_1-z_2)/(z_1-tz_2)$ is a scalar exactly on the hyperbola $st=1$, where it equals $t^{-1}$;
the tempting conclusion is that the modes then generate a quantum torus. \emph{They do not.}
We compute the obstruction exactly: for $s=t^{-1}$ and all $m,n\in\Z$,
\[
  \Hop^{t}_m \Hop^{1/t}_n \;-\; t^{-1}\Hop^{1/t}_n \Hop^{t}_m
  \;=\;\bigl(1-t^{-1}\bigr)\,t^{-m}\, h_{m+n}\bigl[(1+t)X\bigr]\cdot ,
\]
a \emph{multiplication} operator. It vanishes identically if and only if $t=1$, where the
vertex operator degenerates to multiplication by $\sigma[Xz]$ and the statement is trivial.
So the hyperbola carries no quantum-torus structure; what it carries is the statement that the
obstruction is as simple as an operator can be without being zero. Along the way we correct
two errors in the brief that generated this problem.
\end{abstract}

\section{Setup}

Let $\Lam=\Lam_{\mathbb{Q}(t,s)}$ be the ring of symmetric functions over $\mathbb{Q}(t,s)$,
with power sums $p_1,p_2,\dots$, complete homogeneous $h_n$ ($h_n=0$ for $n<0$, $h_0=1$), and
$\sigma[Xz]=\sum_{n\ge0}h_n[X]z^n$. Plethystic notation is as usual: an expression
$f[X+A]$ means substitution $p_k\mapsto p_k+p_k[A]$.

\begin{definition}\label{def:H}
For a parameter $u$, the Hall--Littlewood vertex operator is
\[
  \Hop_u(z)\,f[X]\;=\;\sigma[Xz]\cdot f\!\left[X-\tfrac{1-u}{z}\right],
  \qquad \Hop_u(z)=\sum_{m\in\Z}z^m\,\Hop^{u}_m .
\]
\end{definition}

\begin{lemma}[Modes are well defined]\label{lem:modes}
Let $f\in\Lam$ be of degree $d$. Then $f[X-(1-u)/z]$ is a Laurent polynomial in $z^{-1}$ of
degree at most $d$ with coefficients in $\Lam$, so
\[
  f\!\left[X-\tfrac{1-u}{z}\right]\;=\;\sum_{j=0}^{d}c_j(f)\,z^{-j},
  \qquad
  \Hop^{u}_m f\;=\;\sum_{j=0}^{d}c_j(f)\,h_{m+j}.
\]
In particular $\Hop^{u}_m f=0$ for $m<-d$, and $\Hop^{u}_m f$ is homogeneous of degree $d+m$.
\end{lemma}

\begin{proof}
Write $f$ in the power-sum basis. For a monomial $p_\lambda=\prod_i p_{\lambda_i}$ the
substitution $p_k\mapsto p_k-(1-u^k)z^{-k}$ gives
$\prod_i\bigl(p_{\lambda_i}-(1-u^{\lambda_i})z^{-\lambda_i}\bigr)$, which on expansion over
subsets $S$ of the parts is
$\sum_S \prod_{i\in S}\bigl(-(1-u^{\lambda_i})\bigr)\,p_{\lambda\setminus S}\,z^{-|\lambda_S|}$,
a Laurent polynomial in $z^{-1}$ of degree $\le|\lambda|\le d$. Multiplying by
$\sigma[Xz]=\sum_{n\ge0}h_nz^n$ and extracting $z^m$ gives the stated formula, and the sum is
finite. Since $h_{m+j}=0$ unless $m+j\ge0$ and $j\le d$, we get $\Hop^u_mf=0$ for $m<-d$.
\end{proof}

Consequently, for fixed $f$ the double series $\Hop_t(z_1)\Hop_s(z_2)f$ and
$\Hop_s(z_2)\Hop_t(z_1)f$ are well-defined elements of
$\Lam[[z_1^{\pm},z_2^{\pm}]]$ with each coefficient a \emph{finite} computation:
their supports are the cones
\begin{equation}\label{eq:cones}
  \operatorname{supp}\Hop_t(z_1)\Hop_s(z_2)f\subseteq\{b\ge -d,\ a+b\ge -d\},
  \qquad
  \operatorname{supp}\Hop_s(z_2)\Hop_t(z_1)f\subseteq\{a\ge -d,\ a+b\ge -d\},
\end{equation}
where $(a,b)$ is the exponent of $z_1^az_2^b$. \textbf{The two cones are different.} This is
the whole source of the phenomenon studied in \S3.

We write $\iotain$ for the expansion of a rational function of $z_1,z_2$ as a formal series in
$z_2/z_1$, and $\iotaout$ for the expansion in $z_1/z_2$. Both are ring homomorphisms on the
localisation of $\mathbb{Q}(t,s)[z_1^\pm,z_2^\pm]$ at the relevant denominators.

\section{The two-parameter exchange relation}

\begin{lemma}[Contraction]\label{lem:contract}
For a single letter $v$,
$\ \sigma[-(1-u)v]=\dfrac{1-v}{1-uv}$.
\end{lemma}

\begin{proof}
$\sigma[A]=\exp\bigl(\sum_{k\ge1}p_k[A]/k\bigr)$. For $A=-(1-u)v$ we have $p_k[A]=-(1-u^k)v^k$, so
\[
  \sigma[A]=\exp\Bigl(\sum_{k\ge1}\tfrac{u^kv^k-v^k}{k}\Bigr)
  =\exp\bigl(-\log(1-uv)+\log(1-v)\bigr)=\frac{1-v}{1-uv}. \qedhere
\]
\end{proof}

\begin{remark}\label{rem:brief-inverted}
The brief that generated this problem asserted
$\sigma[-(1-t)z_2/z_1]=\frac{1-tz_2/z_1}{1-z_2/z_1}$, which is the \emph{reciprocal} of the
truth. Followed literally, that sketch produces the relation with the two factors on opposite
sides, i.e.\ $(sz_1-z_2)\Hop_t(z_1)\Hop_s(z_2)=(z_1-tz_2)\Hop_s(z_2)\Hop_t(z_1)$, which is
false (\S4, planted-error control: it fails in $12/12$ tested cases).
\end{remark}

Define the \emph{normal-ordered series}
\begin{equation}\label{eq:Phi}
  \Phi(z_1,z_2)f\;:=\;\sigma[Xz_1]\,\sigma[Xz_2]\;
  f\!\left[X-\tfrac{1-t}{z_1}-\tfrac{1-s}{z_2}\right].
\end{equation}
Note that the argument of $f$ is symmetric under $(z_1,t)\leftrightarrow(z_2,s)$: all of the
asymmetry between the two orderings will be a scalar.

\begin{lemma}[Both orders normal-order to the same $\Phi$]\label{lem:both}
As identities of formal series with coefficients in $\operatorname{End}(\Lam)$,
\[
  \Hop_t(z_1)\Hop_s(z_2)\;=\;\iotain\!\left[\frac{z_1-z_2}{z_1-t z_2}\right]\Phi,
  \qquad
  \Hop_s(z_2)\Hop_t(z_1)\;=\;\iotaout\!\left[\frac{z_1-z_2}{s z_1-z_2}\right]\Phi .
\]
\end{lemma}

\begin{proof}
For the first,
\begin{align*}
  \Hop_t(z_1)\Hop_s(z_2)f
  &=\sigma[Xz_1]\cdot\Bigl(\sigma[Xz_2]f\bigl[X-\tfrac{1-s}{z_2}\bigr]\Bigr)
      \Bigl[X-\tfrac{1-t}{z_1}\Bigr]\\
  &=\sigma[Xz_1]\cdot\sigma\!\left[\Bigl(X-\tfrac{1-t}{z_1}\Bigr)z_2\right]\cdot
      f\!\left[X-\tfrac{1-t}{z_1}-\tfrac{1-s}{z_2}\right]\\
  &=\sigma[Xz_1]\,\sigma[Xz_2]\cdot\sigma\!\left[-\tfrac{(1-t)z_2}{z_1}\right]\cdot
      f\!\left[X-\tfrac{1-t}{z_1}-\tfrac{1-s}{z_2}\right],
\end{align*}
using $\sigma[A+B]=\sigma[A]\sigma[B]$. By Lemma~\ref{lem:contract} with $v=z_2/z_1$ and $u=t$,
the scalar factor is $\frac{1-z_2/z_1}{1-tz_2/z_1}=\frac{z_1-z_2}{z_1-tz_2}$. It arises from
substituting $X\mapsto X-(1-t)/z_1$ into $\sigma[Xz_2]=\sum_{n\ge0}h_n[X]z_2^n$, hence is
produced as a power series in $z_2/z_1$: that is the expansion $\iotain$.

For the second, the same computation with $(z_1,t)\leftrightarrow(z_2,s)$ gives the scalar
$\sigma[-(1-s)z_1/z_2]=\frac{1-z_1/z_2}{1-sz_1/z_2}=\frac{z_1-z_2}{sz_1-z_2}$, produced as a
power series in $z_1/z_2$, i.e.\ $\iotaout$. The function of $f$ and the two $\sigma$'s are
identical in both computations because the argument of $f$ in \eqref{eq:Phi} is symmetric.
\end{proof}

\begin{theorem}[Q99a]\label{thm:A}
For independent parameters $s,t$,
\[
  (z_1-t z_2)\,\Hop_t(z_1)\Hop_s(z_2)
  \;=\;(z_1-z_2)\,\Phi
  \;=\;(s z_1-z_2)\,\Hop_s(z_2)\Hop_t(z_1).
\]
Equivalently, at the level of modes, for all $m,n\in\Z$:
\begin{equation}\label{eq:modeA}
  \Hop^{t}_m \Hop^{s}_n-t\,\Hop^{t}_{m+1}\Hop^{s}_{n-1}
  \;=\;s\,\Hop^{s}_n \Hop^{t}_m-\Hop^{s}_{n-1}\Hop^{t}_{m+1}.
\end{equation}
\end{theorem}

\begin{proof}
Since $\iotain$ is a ring homomorphism and $z_1-tz_2$ is a unit times the denominator,
\[
  (z_1-tz_2)\cdot\iotain\!\left[\frac{z_1-z_2}{z_1-tz_2}\right]
  =\iotain\!\left[(z_1-tz_2)\cdot\frac{z_1-z_2}{z_1-tz_2}\right]=\iotain[z_1-z_2]=z_1-z_2,
\]
and identically $(sz_1-z_2)\cdot\iotaout\bigl[\frac{z_1-z_2}{sz_1-z_2}\bigr]=z_1-z_2$.
Multiply the two identities of Lemma~\ref{lem:both} by $(z_1-tz_2)$ and $(sz_1-z_2)$
respectively. Both give $(z_1-z_2)\Phi$. Extracting the coefficient of $z_1^{m+1}z_2^{n}$
from the outer identity gives \eqref{eq:modeA}; each coefficient is a finite sum by
Lemma~\ref{lem:modes}, so no convergence question arises.
\end{proof}

\begin{corollary}[Zabrocki (2.14)]
Setting $s=t$ gives $(z_1-tz_2)\Hop(z_1)\Hop(z_2)=(tz_1-z_2)\Hop(z_2)\Hop(z_1)$, i.e.
$\Hop_m\Hop_n-t\Hop_{m+1}\Hop_{n-1}=t\Hop_n\Hop_m-\Hop_{n-1}\Hop_{m+1}$.
\end{corollary}

This is the relation recorded in M.~Zabrocki, \emph{On the action of the Hall--Littlewood
vertex operator}, UCSD thesis, \S2.5, eq.~(2.14), read at source; and recorded --- as known
and attributed there to earlier work, not derived independently --- in Korff
\texttt{arXiv:1906.02565} (\texttt{comp.tex} ll.~753--759, ``the latter are \emph{known} to
obey'') and Jing--Liu \texttt{arXiv:2303.10664} (eq.~(e:com1), rearranged, with the $(1-t^n)$
on the creation rather than the annihilation side, attributed to Jing 1991). A fourth
reference carried by the brief, Necoechea--Rozhkovskaya \texttt{arXiv:1902.10049}, does
\emph{not} contain this relation: its only $t$-deformed exchange relation is the
generalized-fermion one for the twisted fields $\Psi^\pm$ on the charged Fock space, which is
inequivalent. Its recovery from
Theorem~\ref{thm:A} is an external consistency check on the whole derivation.

\begin{proposition}[Where the factor is a scalar]\label{prop:hyperbola}
$\dfrac{sz_1-z_2}{z_1-tz_2}$ is independent of $(z_1,z_2)$ if and only if $st=1$, in which
case it equals $t^{-1}$.
\end{proposition}

\begin{proof}
$sz_1-z_2=c(z_1-tz_2)$ for a scalar $c$ forces $c=s$ and $ct=1$, i.e.\ $s=c=t^{-1}$.
Conversely if $s=t^{-1}$ then $sz_1-z_2=t^{-1}(z_1-tz_2)$.
\end{proof}

\section{The hyperbola $st=1$: the defect is not zero}

Fix $s=t^{-1}$, $t\neq0$. Proposition~\ref{prop:hyperbola} and Theorem~\ref{thm:A} give
\[
  (z_1-tz_2)\Bigl(\Hop_t(z_1)\Hop_{1/t}(z_2)-t^{-1}\Hop_{1/t}(z_2)\Hop_t(z_1)\Bigr)=0,
\]
and it is tempting to divide by $(z_1-tz_2)$ and conclude the quantum-torus relation
$\Hop^t_m\Hop^{1/t}_n=t^{-1}\Hop^{1/t}_n\Hop^t_m$. \textbf{This division is illegitimate},
because by \eqref{eq:cones} the difference is supported in a region on which
$(z_1-tz_2)$ is a zero divisor: the series $\delta(tz_2/z_1)=\sum_{j\in\Z}t^jz_1^{-j}z_2^{j}$
satisfies $(z_1-tz_2)\delta(tz_2/z_1)=0$. We now compute the defect exactly.

\begin{notation}
For $N\ge0$ set $G_N:=\sum_{a+b=N,\;a,b\ge0}t^a h_a h_b$, and $G_N:=0$ for $N<0$. Equivalently
$G_N=h_N[(1+t)X]$, where $(1+t)X$ denotes the alphabet $X\cup tX$, i.e.\ the plethystic
substitution $p_k\mapsto(1+t^k)p_k$; indeed
$\sum_N G_Nz^N=\sigma[tXz]\sigma[Xz]=\sigma[(1+t)Xz]$.
\end{notation}

\begin{lemma}[The $\delta$-term]\label{lem:delta}
$\displaystyle
 \bigl(\iotain-\iotaout\bigr)\!\left[\frac{z_1-z_2}{z_1-t z_2}\right]
 =\bigl(1-t^{-1}\bigr)\sum_{j\in\Z}t^{j}z_1^{-j}z_2^{j}.
$
\end{lemma}

\begin{proof}
$\iotain\bigl[\tfrac{1}{z_1-tz_2}\bigr]=\sum_{k\ge0}t^kz_2^kz_1^{-k-1}$ and
$\iotaout\bigl[\tfrac{1}{z_1-tz_2}\bigr]=-\sum_{k\ge0}t^{-k-1}z_1^kz_2^{-k-1}$; their
difference is $\sum_{j\in\Z}t^jz_1^{-j-1}z_2^j$ (the terms $j\ge0$ from the first sum, the
terms $j=-k-1<0$ from the second). Multiplying by the polynomial $z_1-z_2$,
\[
  \sum_{j\in\Z}t^jz_1^{-j}z_2^j-\sum_{j\in\Z}t^jz_1^{-j-1}z_2^{j+1}
  =\sum_{j\in\Z}\bigl(t^j-t^{j-1}\bigr)z_1^{-j}z_2^{j},
\]
after reindexing $j\mapsto j-1$ in the second sum. \qedhere
\end{proof}

\begin{lemma}[Diagonal cancellation]\label{lem:diag}
Let $s=t^{-1}$ and let $\Phi_{a,b}$ denote the coefficient of $z_1^az_2^b$ in $\Phi$. Then for
every $N\in\Z$ and every $f\in\Lam$, the sum $\sum_{l\in\Z}t^{l}\Phi_{l,N-l}f$ is finite and
\[
  \sum_{l\in\Z}t^{l}\,\Phi_{l,\,N-l}\,f\;=\;G_N\cdot f .
\]
\end{lemma}

\begin{proof}
Expand $f$ in the power-sum basis and apply the two-variable analogue of the expansion in
Lemma~\ref{lem:modes}: assigning each part of $\lambda$ to ``keep'', ``$z_1$'' or ``$z_2$'',
\[
  f\!\left[X-\tfrac{1-t}{z_1}-\tfrac{1-s}{z_2}\right]
  =\sum_{(\mu,i,i')}c_{\mu,i,i'}\;p_\mu\,z_1^{-i}z_2^{-i'},
  \qquad i,i'\ge0,\ i+i'\le d .
\]
Hence $\Phi_{a,b}f=\sum c_{\mu,i,i'}\,p_\mu\,h_{a+i}h_{b+i'}$, which vanishes unless
$a\ge-d$ and $b\ge-d$; so the sum over $l$ is finite. Substituting $a=l+i$,
\[
  \sum_{l}t^l\Phi_{l,N-l}f
  =\sum_{(\mu,i,i')}c_{\mu,i,i'}\,t^{-i}p_\mu\sum_{a\ge0}t^{a}h_ah_{N+i+i'-a}
  =\sum_{w\ge0}\Bigl(\sum_{i+i'=w}c_{\mu,i,i'}t^{-i}\Bigr)p_\mu\,G_{N+w}.
\]
Now the inner bracket is exactly the coefficient of $z_2^{-w}$ in the specialisation
$z_1\mapsto tz_2$ of the displayed expansion. For $s=t^{-1}$ that specialisation is
\[
  -\frac{1-t^k}{(tz_2)^k}-\frac{1-t^{-k}}{z_2^{k}}
  =\frac{-t^{-k}+1-1+t^{-k}}{z_2^{k}}=0
  \qquad\text{for every }k\ge1,
\]
i.e.\ $f[X-\tfrac{1-t}{z_1}-\tfrac{1-s}{z_2}]\big|_{z_1=tz_2}=f[X]$. Therefore the bracket
vanishes for $w>0$ and equals the coefficient $f_\mu$ of $p_\mu$ in $f$ for $w=0$. Only the
$w=0$ term survives, giving $\sum_\mu f_\mu p_\mu\,G_N=G_N\cdot f$.
\end{proof}

\begin{theorem}[Q99b]\label{thm:B}
Let $s=t^{-1}$ with $t\neq0$. Then for all $m,n\in\Z$, as operators on $\Lam$,
\[
  \boxed{\;
  \Hop^{t}_m \Hop^{1/t}_n \;-\; t^{-1}\,\Hop^{1/t}_n \Hop^{t}_m
  \;=\;\bigl(1-t^{-1}\bigr)\,t^{-m}\;h_{m+n}\bigl[(1+t)X\bigr]\cdot\;}
\]
where the right-hand side is the operator of multiplication by the symmetric function
$h_{m+n}[(1+t)X]=G_{m+n}$ (which is $0$ when $m+n<0$).
Equivalently, in generating-function form,
\[
  \Hop_t(z_1)\Hop_{1/t}(z_2)-t^{-1}\Hop_{1/t}(z_2)\Hop_t(z_1)
  =\bigl(1-t^{-1}\bigr)\,\delta\!\left(\tfrac{tz_2}{z_1}\right)\,\sigma\bigl[(1+t)Xz_2\bigr]\cdot .
\]
\end{theorem}

\begin{proof}
By Lemma~\ref{lem:both} with $s=t^{-1}$ we have
$\frac{z_1-z_2}{sz_1-z_2}=\frac{t(z_1-z_2)}{z_1-tz_2}$, hence
$t^{-1}\Hop_{1/t}(z_2)\Hop_t(z_1)=\iotaout\bigl[\frac{z_1-z_2}{z_1-tz_2}\bigr]\Phi$, and so
\[
  C:=\Hop_t(z_1)\Hop_{1/t}(z_2)-t^{-1}\Hop_{1/t}(z_2)\Hop_t(z_1)
  =\bigl(\iotain-\iotaout\bigr)\!\left[\frac{z_1-z_2}{z_1-tz_2}\right]\Phi .
\]
By Lemma~\ref{lem:delta} this equals $(1-t^{-1})\sum_{j\in\Z}t^jz_1^{-j}z_2^{j}\,\Phi$.
Extract the coefficient of $z_1^mz_2^n$: the pairs contributing are $(a,b)=(m+j,\,n-j)$, so
\[
  C_{m,n}f=\bigl(1-t^{-1}\bigr)\sum_{j\in\Z}t^{j}\,\Phi_{m+j,\,n-j}f ,
\]
a finite sum because $\Phi_{a,b}f=0$ unless $a\ge-d$ and $b\ge-d$ (Lemma~\ref{lem:diag}),
which confines $j$ to $-d-m\le j\le n+d$. Writing $j=l-m$ and applying Lemma~\ref{lem:diag}
with $N=m+n$,
\[
  C_{m,n}f=\bigl(1-t^{-1}\bigr)\,t^{-m}\sum_{l\in\Z}t^{l}\Phi_{l,\,m+n-l}f
  =\bigl(1-t^{-1}\bigr)\,t^{-m}\,G_{m+n}\cdot f . \qedhere
\]
\end{proof}

\begin{corollary}[The quantum torus does \emph{not} appear]\label{cor:notorus}
On the hyperbola $st=1$ the relation
$\Hop^t_m\Hop^{1/t}_n=t^{-1}\Hop^{1/t}_n\Hop^t_m$ holds for all $m,n$ if and only if $t=1$.
\end{corollary}

\begin{proof}
By Theorem~\ref{thm:B} the defect is $(1-t^{-1})t^{-m}G_{m+n}$. Since $G_0=1$, taking
$m+n=0$ gives defect $(1-t^{-1})t^{-m}\neq0$ unless $t=1$. Conversely at $t=1$ the prefactor
$1-t^{-1}$ vanishes, so the defect is identically zero.
\end{proof}

\begin{remark}[$t=1$ is degenerate, and the brief's negative control was wrong]
\label{rem:control}
The brief asserted: ``At $t=s=1$ the hyperbola is met, $1/t=1$, and the naive corollary
predicts that the Bernstein operators commute. \emph{They do not} \dots\ Any derivation of the
mode relation that does not break at $t=1$ is wrong.''

This is false, and following it would have caused me to reject the correct theorem. At $t=1$
Definition~\ref{def:H} degenerates: $\Hop_1(z)f=\sigma[Xz]\,f[X-0]=\sigma[Xz]f$, so
$\Hop^{1}_m$ is \emph{multiplication by $h_m$} --- verified independently in \S4 --- and
multiplication operators commute. The brief's supporting observation, that (2.14) at $t=1$ is
strictly weaker than commutativity, is true but does not entail non-commutativity:
commutativity at $t=1$ is an extra true fact that (2.14) alone does not see. Corollary
\ref{cor:notorus} is therefore consistent with the $t=1$ specialisation, and the prediction
does move: the prefactor $1-t^{-1}$ takes the values $0,\ \tfrac12,\ 2,\ -1$ at
$t=1,2,-1,\tfrac12$.
\end{remark}

\begin{remark}[What the hyperbola actually carries]
The positive content of Theorem~\ref{thm:B} is not a relation but a \emph{regularity}. For
generic $(s,t)$ the two orderings differ by a $(z_1,z_2)$-dependent factor and no mode-level
statement of quantum-torus type can even be formulated. On $st=1$ one can formulate it, and it
fails --- but it fails by a multiplication operator, the simplest kind of endomorphism of
$\Lam$ there is. The mechanism is exactly the plethystic cancellation in
Lemma~\ref{lem:diag}: on the diagonal $z_1=tz_2$, which is where the $\delta$-term is
supported, the two shifts $-\frac{1-t}{z_1}$ and $-\frac{1-t^{-1}}{z_2}$ annihilate each
other, so $f$ passes through untouched and only the two $\sigma$'s survive. The hyperbola is
precisely the locus where the singular support of the exchange and the vanishing locus of the
total plethystic shift coincide.
\end{remark}

\begin{remark}[$t=-1$, the free-fermion point]\label{rem:fermion}
At $t=-1$ (the intersection of $st=1$ with the diagonal $s=t$, other than $t=1$) the
substitution $p_k\mapsto(1+(-1)^k)p_k$ kills all odd power sums, so $G_N=0$ for $N$ odd and
$G_{2M}=h_M[X^2]$ (the alphabet of squares). The defect is therefore
$2\,(-1)^{m}\,h_{(m+n)/2}[X^2]$ for $m+n$ even and $0$ for $m+n$ odd. For $m=n=1$ this is
$-2p_2$, which is consistent with the known failure $\Hop_1\Hop_1\cdot 1=-p_2\neq0$ of the
naive anticommutation reading at $t=-1$.
\end{remark}

\section{Verification}

All computations were done in an engine written for this session
(\texttt{projects/scratch/q99/engine.py}): symmetric functions as dictionaries
\{partition $\to$ coefficient\} in the power-sum basis, $h_n=\sum_{\lambda\vdash n}p_\lambda/z_\lambda$,
and the plethystic shift implemented by explicit subset expansion. It shares no code path with
the pre-existing scripts \texttt{/tmp/browse/0907pm/hlvo.py} and \texttt{hlvo2.py} (which use
\texttt{sympy.subs}, \texttt{Poly}, and $\exp(\sum p_ku^k/k)$); those were read but not run,
and nothing here depends on them.

\medskip\noindent\textbf{Engine anchors} (untuned; neither was used to build the engine):
\begin{itemize}
\item $t=1$: $\Hop^{1}_mf=h_mf$ for $m\in\{0,1,2,3\}$ and $f\in\{1,p_1,p_1^2,p_2\}$ --- $16/16$.
\item $t=0$: the classical Bernstein operator, $\Hop^{0}_m s_\lambda=s_{(m,\lambda)}$ for
$m\ge\lambda_1$, with $s_\lambda$ computed independently by Jacobi--Trudi --- $13/13$;
straightening ($m<\lambda_1$) also behaves correctly, e.g.\ $\Hop^0_1s_{(2)}=0$ and
$\Hop^0_0s_{(2)}=-s_{(1,1)}$.
\end{itemize}

\medskip\noindent\textbf{Theorem~\ref{thm:A}} (mode form \eqref{eq:modeA}), generic symbolic
$t,s$: $126$ cases $=6$ arguments $f\in\{1,p_1,p_2,p_1^2,p_3,p_2p_1\}$ $\times$ $21$ pairs
$(m,n)$ including negative modes. Residual identically $0$ in all $126$.

\smallskip\noindent\emph{Non-degeneracy.} At $(m,n)=(1,2)$, $f=p_1$ the four terms of
\eqref{eq:modeA} are pairwise distinct and each depends on both $t$ and $s$; the bare
commutator $\Hop^t_1\Hop^s_2f-\Hop^s_2\Hop^t_1f$ is nonzero. The test is not vacuous.

\smallskip\noindent\emph{Planted-error control.} Replacing the relation by
$(z_1-\alpha z_2)\Hop_t\Hop_s=(\beta z_1-\gamma z_2)\Hop_s\Hop_t$ for
$(\alpha,\beta,\gamma)\in\{(s,t,1),(t,t,1),(t,1/s,1),(ts,s,1),(t,s,s)\}$ gives a nonzero
residual in $12/12$ tested cases for each of the five perturbations; only $(t,s,1)$ passes.
In particular the inverted sketch of Remark~\ref{rem:brief-inverted} --- which is the variant
$(s,t,1)$ --- fails $12/12$.

\medskip\noindent\textbf{Lemma~\ref{lem:both}} (the proof, not just the statement): both
$\Hop^t_m\Hop^s_nf-t\Hop^t_{m+1}\Hop^s_{n-1}f$ and
$s\Hop^s_n\Hop^t_mf-\Hop^s_{n-1}\Hop^t_{m+1}f$ were checked against
$\Phi_{m,n}f-\Phi_{m+1,n-1}f$ with $\Phi$ computed directly from \eqref{eq:Phi}:
$65/65$ agree for each ordering ($5$ arguments $\times$ $13$ mode pairs).

\medskip\noindent\textbf{Theorem~\ref{thm:B}}, generic symbolic $t$ with $s=t^{-1}$:
$138$ cases $=6$ arguments $\times$ $23$ pairs $(m,n)$, mismatches $0$. The defect was
\emph{nonzero} in $120$ of the $138$; the $18$ vanishing cases factor as $3\times 6$ --- the
three pairs with $m+n<0$, times the six arguments --- exactly as Theorem~\ref{thm:B} predicts,
and for no other reason. A higher-degree stress test (arguments $p_2^2,p_4,p_2p_1^2,p_3p_2$
of degrees $4$ and $5$; modes up to $5$ and down to $-2$) gives a further $28/28$.

\medskip\noindent\textbf{The two intermediate steps of the proof of Theorem~\ref{thm:B}
were checked separately}, so that the final agreement is not the only evidence:
$C_{m,n}=(1-t^{-1})\sum_jt^j\Phi_{m+j,n-j}$ ($21/21$) and Lemma~\ref{lem:diag}
$\sum_lt^l\Phi_{l,N-l}=G_N\cdot$ ($15/15$). The two expressions for $G_N$
($\sum_{a+b=N}t^ah_ah_b$ and $h_N[(1+t)X]$) agree for $0\le N\le5$.

\medskip\noindent\textbf{Controls on the hyperbola} (Remark~\ref{rem:control}): at $t=1$ the
defect is identically $0$ over all tested $(m,n)$ and $f$; at $t=2$, $t=-1$ and $t=\frac12$ it
is nonzero, and matches Theorem~\ref{thm:B} in every case. At $t=-1$ it vanishes exactly on
$m+n$ odd, as predicted.

\section{Relation to the ribbon-side $ts=1$ locus, and what is \emph{not} claimed}

PROVE 2026-09-07~c2 (\texttt{proofs/2026-09-07-c2-Q96-order-over-sublattice.tex},
Thm~4.1(2), Cor~4.2(iii)) established, on two-bead matrix elements of the ribbon
operators, that
$\langle M'|[R_e(t),R_f(s)]|M\rangle=\sum_{(b,c)}t^{P-k}s^{Q}(1-(ts)^k)$, hence that the
\emph{two-bead sector} vanishes identically iff $ts=1$. It is only that sector: the one-bead
sector of the same commutator is nonzero at every point of $ts=1$ except $t=-1$ --- verified
in the review of 8 September 2026 on $1232/1232$ entries ($|\lambda|\le6$, $1\le e,f\le4$),
where it is moreover divisible by $(1+t)$ on all of them. The hyperbola is thus named by two
computations. Are they two witnesses or one witness twice?

\noindent\textbf{They are not the same witness}, and this session's result is what
distinguishes them:
\begin{itemize}
\item On the ribbon side the vanishing object is one \emph{sector} of the commutator, not the
commutator itself; and the two sides' degenerate points along the hyperbola are exactly
swapped. On $ts=1$ the ribbon commutator $[R_e(t),R_f(1/t)]$ is nonzero at $t=1$ ($400$
entries, $|\lambda|\le5$, $1\le e,f\le4$) and identically zero at $t=-1$; the twisted
vertex-operator commutator is zero at $t=1$ (Corollary~\ref{cor:notorus}) and nonzero at
$t=-1$ ($-2p_2$ at $m=n=1$, Remark~\ref{rem:fermion}). Complementary vanishing loci on the
same hyperbola: no reparametrisation carries one derivation to the other.
\item On the vertex-operator side the plain commutator does not vanish on $ts=1$: e.g.
$[\Hop^t_1,\Hop^{1/t}_1]\cdot1=\frac{t^2-1}{t}h_2$, which is nonzero for $t\neq\pm1$, and
$[\Hop^t_2,\Hop^{1/t}_1]\cdot1\neq0$ at $t=-1$. The natural $t^{-1}$-twisted commutator does
not vanish on $ts=1$ either, by Corollary~\ref{cor:notorus}: it vanishes only at the single
point $t=1$.
\end{itemize}
So the two occurrences of $st=1$ have different content. The vertex-operator occurrence is
Proposition~\ref{prop:hyperbola} --- a statement about when a rational exchange factor
degenerates to a scalar --- and the mode-level consequence is a \emph{failure} with a
computed defect, not a vanishing.

\medskip\noindent\textbf{Explicitly not claimed here.} That $st=1$ is ``the hidden structure
behind the ribbon commutator''; that $R_e(t)$ is a mode $\Hop^t_m$ or any specialisation of
one; that Theorem~\ref{thm:B} transports to the ribbon side. $R_e(t)$ is the connected
truncation of the ribbon-move expansion and the Q92/Q96 enumeration runs over single bead
moves, not over $\sigma[Xz]$; a transport would be a computation, and it was not done.
Theorem~\ref{thm:A} warrants Theorem~\ref{thm:A}.

\section{Gaps and what was not attempted}

\begin{enumerate}
\item \textbf{Not attempted: the transport to $R_e(t)$.} See \S5. This is the natural next
question and the obstruction is named there.
\item \textbf{Not attempted: a basis-level interpretation.} Theorem~\ref{thm:B} says the
defect is multiplication by $h_{m+n}[(1+t)X]$. The alphabet $(1+t)X$ is the one for which
$\prod_i(1-x_iz)^{-1}(1-tx_iz)^{-1}$ is the generating function; it is natural to ask whether
the defect is the image of a Hall--Littlewood-theoretic operator (a Macdonald $q,t$
specialisation, or the $t$-analogue of a Boson--Fermion normal-ordering constant). Not
investigated.
\item \textbf{Coverage.} Theorem~\ref{thm:A} and Theorem~\ref{thm:B} are proved for all
$m,n\in\Z$ and all $f\in\Lam$; the computations are checks, not the proof. The checks cover
arguments of degree $\le5$ and modes in $[-2,5]$. No claim is made about specialisations where
$t=0$ (Definition~\ref{def:H} still makes sense, but $s=t^{-1}$ does not).
\item \textbf{The $\iota$ formalism.} Lemma~\ref{lem:both} identifies which expansion region
each ordering produces. I have argued this from the derivation (the factor arises by
substituting into $\sigma[Xz_2]=\sum_{n\ge0}h_nz_2^n$, a power series in $z_2$); the support
statement \eqref{eq:cones}, proved from Lemma~\ref{lem:modes}, is the independent confirmation,
and the $65/65$ numerical agreement in \S4 is a third. I regard this as closed but flag it as
the one step where an expansion-region error would be invisible in the final formula.
\end{enumerate}

\end{document}
